% related_final.tex — loaded by main_final.tex

\paragraph{Chain-of-Thought Reasoning.}
Chain-of-thought (CoT) prompting~\citep{wei2022chain} and its variants~\citep{wang2023selfconsistency, yao2023tree} have become the de facto approach for improving LLM reasoning.
Recent work has moved beyond prompting to \emph{training} models with explicit reasoning traces, yielding ``thinking'' LLMs such as DeepSeek-R1~\citep{deepseek2025r1}, QwQ~\citep{qwen2025qwq}, and Qwen3~\citep{qwen2025qwen3}.
These models produce a \texttt{<think>...</think>} block before the final answer.
While effective at unconstrained budgets, we show that this thinking overhead imposes a severe penalty under realistic token constraints.

\paragraph{Test-Time Compute Scaling.}
The scaling of inference-time computation has emerged as a complementary axis to model scaling~\citep{snell2024scaling}.
Best-of-N sampling~\citep{lightman2024lets}, process reward models~\citep{wang2024math}, and iterative refinement~\citep{madaan2023selfrefine} all trade additional tokens for improved accuracy.
However, these approaches assume that \emph{more tokens always help}.
Our findings challenge this assumption: at constrained budgets, the overhead of structured reasoning can actively \emph{hurt} performance.

\paragraph{Adaptive Computation and Early Exit.}
Adaptive computation~\citep{graves2016adaptive, dehghani2018universal} allows models to allocate variable compute per input.
In the LLM setting, speculative decoding~\citep{leviathan2023fast, chen2023accelerating} and early exit mechanisms~\citep{schuster2022confident, bae2023fast} reduce inference cost.
Our work differs in that we identify a \emph{free} confidence signal---natural stopping---that requires no additional training or auxiliary models, and propose routing between thinking and non-thinking modes rather than adjusting depth within a single forward pass.

\paragraph{Budget-Aware Reasoning.}
A growing body of concurrent work addresses reasoning efficiency.
AdaThink~\citep{adathink2025} trains budget controllers to allocate thinking tokens adaptively; Think Less, Think Better~\citep{hou2025thinkless} injects wrap-up tokens to shorten chains; Satisficing~\citep{luo2025satisficing} applies confidence-based early exit; and SelfBudgeter~\citep{liu2025selfbudgeter} learns per-instance budgets.
These methods all operate \emph{within} the thinking paradigm---they assume thinking mode is desirable and seek to make it cheaper.
Our work challenges the premise itself: we show that at practical budgets, the \emph{entire thinking framework} is counterproductive, and propose routing \emph{between} thinking and non-thinking modes rather than optimizing within thinking mode alone.

\paragraph{Overthinking and Compute Waste.}
\citet{chen2024overthinking} identify overthinking in mathematical reasoning, where models produce redundant reasoning steps.
\citet{nvidia2025pareto} propose a Pareto framework for allocating test-time compute.
Our contribution is orthogonal and more fundamental: rather than reducing waste within thinking mode, we show that \emph{disabling} thinking mode entirely yields dramatic improvements at constrained budgets.
The natural-stop oracle we identify provides a free, training-free signal for routing decisions---a capability not explored in prior work.
